\textbf{Frente al estado del arte.} Los modelos espaciotemporales de la Sección~\ref{sec:arte} se evalúan sobre autopistas con datos reales de detectores, como PeMS o METR-LA, y reportan su error como un número aislado. Nuestro MAE no es comparable en valor absoluto con esos trabajos: la red es urbana, los datos son sintéticos y el horizonte es fijo a 60~min. La diferencia está en el enfoque. Predecir la corrección sobre el promedio histórico permite un modelo de 29\,908 parámetros que gana justo donde lo habitual falla (congestión y lluvia) y que se mantiene con solo una parte de los vehículos conectados. En detección, la comparación con un detector al azar con la misma FAR es una práctica que rara vez se reporta, y muestra que una DR alta puede deberse solo a disparar muchas alarmas. En el aviso, a diferencia de los estudios que evalúan el rerouting con información perfecta, aquí la acción depende de un detector imperfecto, y su costo, los avisos por falsas alarmas, queda incluido en el resultado.

\textbf{Compromiso entre detectar y no molestar.} Bajar la FAR de 6,5 a 4,4~\% en los días cargados retrasó la detección en los días tranquilos, donde la cola se forma despacio (MTTD de 7,3~min). En el SIL se usó un umbral más estricto, y aun así el 69~\% de los avisos fueron por falsas alarmas. Que la red mejore de todos modos se debe a que los avisos caducan y el navegador solo desvía si se gana tiempo. Para un sistema operable hay que reducir esas falsas alarmas.

\textbf{Validez del dato sintético.} El gemelo digital reproduce órdenes de magnitud de Cuenca (Anexo~\ref{an:valida}), pero no su variabilidad real: no hay buses, motos, eventos especiales, fallas de sensores ni conductores que se desvíen por su cuenta. Los incidentes son vehículos detenidos con duración conocida. Por eso los resultados se presentan como sintéticos, y la cifra relevante es la mejora relativa frente a la línea base en las mismas condiciones, no el valor absoluto.

\textbf{Alcance del bucle cerrado.} El SIL evalúa un solo cierre total con una semilla. Demuestra que el mecanismo funciona, pero no basta para estimar la mejora media sobre muchos incidentes; ese es el siguiente experimento, y el script ya acepta cualquier incidente (\texttt{--inc}).
